\documentclass[12pt, a4paper]{article}
\usepackage{graphicx} % Required for inserting images
\graphicspath{{images/}{images-desafios/}}
\usepackage[a4paper, total={7.5in, 10in}]{geometry}
\usepackage{float} % pra conseguir usar o \begin{figure}[H]
\usepackage{indentfirst} % pra fazer tab n 1 paragrafos dos capitulos
% package pra tabela
\usepackage{multirow}
\usepackage{array}

%pacote pra codigo
\usepackage{listings} 
\usepackage{xcolor}
\definecolor{codegreen}{rgb}{0,0.6,0}
\definecolor{codegray}{rgb}{0.5,0.5,0.5}
\definecolor{codepurple}{rgb}{0.58,0,0.82}
\definecolor{backcolour}{rgb}{0.95,0.95,0.92}

\lstdefinestyle{myCodeStyle}{
	backgroundcolor=\color{backcolour},   
	commentstyle=\color{codegreen},
	keywordstyle=\color{cyan},
	numberstyle=\tiny\color{codegray},
	stringstyle=\color{codepurple},
	basicstyle=\ttfamily\footnotesize,
	breakatwhitespace=false,         
	breaklines=true,                 
	captionpos=b,                    
	keepspaces=true,                 
	%numbers=left,                    
	numbersep=5pt,                  
	showspaces=false,                
	showstringspaces=false,
	showtabs=false,                  
	tabsize=2
}
\lstset{style=myCodeStyle}
\renewcommand{\lstlistingname}{Código}%Muda o termo "Listing":  Listing -> Código
\renewcommand{\lstlistlistingname}{Lista de \lstlistingname s}% List of Listings -> Lista de Códigos


\newcommand{\unit}[1]{\ensuremath{\, \mathrm{#1}}} % comando pra tirar do math mode a unidade de medida

\def\tamanhoGrafico{0.8}

\newcommand{\ohm}{\unit{\Omega} }
\newcommand{\kohm}{\unit{k\Omega}}

\title{
	Laboratório Digital A - PCS3335
	\\ Relatório da Experiência 4
}

\author{
	Alunos:
	\\ Rafael Hipólito Rodrigues - N.USP: 14604308
	\\ {{PERSON-2}} - N.USP: {{PERSON-ID-2}}
	\\ \\
	Turma 3 - Bancada A6
	\\ Professor: Bruno de Carvalho Albertini
}

%\date{}

\begin{document}
	
	\maketitle
	
	\section*{Planejamento} 
	Elaboração de uma UART simplificada. A UART foi dividida em \texttt{Baud Rate Generator  (BRG)}, 
	\texttt{Transmitter Timing \& Control (TTC)},  e \texttt{Transmitter Shift Register (TSR)}. O Baud Rate Generator divide
	o clock de 1.8432MHhz em 12 vezes. Os dados são envidados em 9600bps. 
	\par
	\texttt{Desafio:} Iniciar a transmição dos bits ao comandado de uma chave e ao final da transmissão acender um LED.
	
	
	\section*{Diagramas}
	
	\subsection*{FSM}
	
	Diagrama da Maquina de Estados Finitos gerado no Quartus: 
	
	\begin{figure}[H]
		\centering
		\includegraphics[width=0.7\linewidth]{images/FSM_exp4.png}
		\caption{FSM da uart.}
		\label{fig:FSM_exp4}
	\end{figure}
	
	
	
	\subsection*{RTL}
	
	Diagramas RTL da entidades gerado no Quartus:
	
	\begin{figure}[H]
		\centering
		\includegraphics[width=0.7\linewidth]{images/RTL_uart.png}
		\caption{Diagrama RTL da UART.}
		\label{fig:RTL_uart}
	\end{figure}
	
	\begin{figure}[H]
		\centering
		\includegraphics[width=0.7\linewidth]{images/RTL_BRG.png}
		\caption{Diagrama RTL do Baud Rate Generator.}
		\label{fig:RTL_BRG}
	\end{figure}
	
	\begin{figure}[H]
		\centering
		\includegraphics[width=0.7\linewidth]{images/RTL_transmitter.png}
		\caption{Diagrama RTL do Transmissor Serial.}
		\label{fig:RTL_transmitter}
	\end{figure}
	
	\begin{figure}[H]
		\centering
		\includegraphics[width=0.7\linewidth]{images/RLT_TTC.png}
		\caption{Diagrama RTL do Transmitter Timing \& Control - TCC.}
		\label{fig:RLT_TTC}
	\end{figure}
	
	Os contadores e registradores utilizados foram os mesmos da experiência 2.
	
	\section*{Pinagem}

	
	\begin{table}[H]
		\begin{center}
			\begin{tabular}{|l|l|l|}
				\hline
				clock       & PIN\_M9  & Input  \\ \hline
				go          & PIN\_U13 & Input  \\ \hline
				readyLed    & PIN\_AA2 & Output \\ \hline
				reset       & PIN\_B16 & Input  \\ \hline
				serial\_out & PIN\_N16 & Output \\ \hline
			\end{tabular}
		\end{center}
	\end{table}
	
	
	
	\section*{Testes}

	\begin{itemize}
		\item bits adequadamente envidados em 9600bps.
		\par Resultado: SUCESSO.
		
		\item bits de "10010000100" envidados corretamente.
		\par Resultado: SUCESSO.
		
		\item Ao ativar e desativa o reset, os dados são enviados novamente. 
		\par Resultado: SUCESSO.
		
	\end{itemize}
	
	\subsection*{Comentários}
	
	Os resultados de todos os testes foram de acordo com o planejado.
	
\end{document}
